
% ============================================================================
\clearpage
\section*{v4.9.0-r83 Source-Only: Quantum Decay of the Matter--Tachyon Vacuum, the Tachyon Image of an Action, and a Correction to r82}
\addcontentsline{toc}{section}{v4.9.0-r83 Source-Only: Quantum Decay of the Matter--Tachyon Vacuum, the Tachyon Image of an Action, and a Correction to r82}
\label{sec:v490-r83-vacuum-decay-image}
% ============================================================================

\paragraph{Purpose.}
Layer v4.9.0-r82 left the quantum decay of the vacuum of the matter--tachyon pairing open
(\S\ref{sec:v490-r82-higgs-stability}).  On 5 November 2026 the author chose to check it first.  The check (H76) needed a preferred
frame, and the search for one exposed that the tachyon image of an \emph{action} had not been fixed.  The author asked to settle this
before choosing.  This layer records both pre-registered tests (H76, H77, both not blind), an exploration of ghost condensation as the
frame-defining mechanism, and a correction to r82.1(b), (c).  The notes are \texttt{notes/\allowbreak research/\allowbreak v4\_10\_0\_r83/}.
Earlier layers are not rewritten.

\noindent\textbf{Source-only status.}
This v4.9.0-r83 layer is source-only.  It does not build a PDF, does not assign a new DOI,
and does not authorize a public release.  The v4.9.0 Version DOI
\href{https://doi.org/10.5281/zenodo.23534954}{10.5281/zenodo.23534954} and the Concept DOI
\href{https://doi.org/10.5281/zenodo.19210728}{10.5281/zenodo.19210728} are unchanged by it.

\subsection*{v4.9.0-r83 results}
\label{subsec:v490-r83-results}

\paragraph{Theoremlet v4.9.0-r83.1 (H76: gravitational decay of the vacuum).}
Pre-registered (commit \texttt{2b47e7f}, not blind; the scaling is known from Cline, Jeon and Moore, hep-ph/0311312).  One massless
scalar in M (positive energy) and one in T (negative energy), coupled only by a graviton, \(\kappa^2=32\pi/M_P^2\); the process is
\(0\to\varphi\varphi\,\chi\chi\).
\begin{enumerate}[leftmargin=2em]
\item[(a)] With only a Lorentz-invariant cutoff \(P^2<\Lambda^2\) the decay is allowed and its rate is unbounded: with an additional frame
cutoff \(L\Lambda\) the ratios \(\Gamma(L)/\Gamma(L/2)\) are \(7.06,4.77,4.23,4.07,4.02,4.006,4.002\) for \(L=2,\dots,128\).  The rate grows as
\(L^2\), the volume of boosts on the hyperboloid.
\item[(b)] With a cutoff \(|k|<\Lambda\) in a preferred frame, \(\Gamma/V=c\,\Lambda^8/M_P^4\) with exponent \(8.000\) and \(7.997\) and
\(c=2.25\times10^{-5}\) (relative error \(0.12\%\)).
\item[(c)] Transferred to the gamma-ray bound of Cline et al.\ (their \(c\sim1\), \(\Lambda\lesssim3\) MeV) the bound moves by
\(c^{-1/7}=4.6\) to \(\Lambda\lesssim14\) MeV (rough).  A T sector that is a Lorentz-invariant copy up to the Higgs scale and couples only
gravitationally contradicts the gamma-ray background.
\end{enumerate}
Run~1 (direct Monte Carlo) was precise for (b) but not for (a) at large \(L\); run~2 with ten times more samples did not converge
(heavy-tailed weights).  Run~3 computes the same quantity by Lorentz invariance (configuration in the rest frame of the matter pair,
exact boost volume) and agrees with run~1 at \(L=1,2,4\) to better than \(1\%\).  The criterion was not changed; the method was changed
after seeing the data.

\paragraph{Theoremlet v4.9.0-r83.2 (the condensate of I4 defines no frame).}
Exploration.  A condensate \(\langle\chi\rangle=\text{const}\) of tachyonic modes is Lorentz invariant.  It selects no frame, and the
decay of r83.1 persists.  It can supply a scale, not the frame.

\paragraph{Theoremlet v4.9.0-r83.3 (H77: the tachyon image of an action).}
Pre-registered (commit \texttt{d087d67}, not blind).  Three readings for a scalar with \(L_M=\tfrac12(\chi_t^2-\chi_x^2)-\tfrac12m^2\chi^2\):
A, \(L_T=-L_M\) (energy parity, used in r82.1(b)); B, the kinematics continued (\(X\mapsto-X\), constants and potentials kept), as the
rapidity \(\eta+i\pi/2\) of r20.2 (v4.8.0) gives \(p\mapsto ip\); B\('\), \(m^2\mapsto-m^2\) with real kinematics.  Criteria: K1 tachyonic
dispersion (r07.2, r20.2, H71); K2 massless \(L_T=-L_M\) (r81); K3 a running mode carries energy of the opposite sign (F3 on readouts).
\begin{center}
\small
\begin{tabular}{@{}lcccc@{}}
\toprule
reading & dispersion & K1 & K2 & K3\\
\midrule
A & \(\omega^2=k^2+m^2\) & no & yes & yes\\
B & \(\omega^2=k^2-m^2\) & yes & yes & yes\\
B\('\) & \(\omega^2=k^2-m^2\) & yes & no & no\\
\bottomrule
\end{tabular}
\end{center}
Only reading B is compatible.  Reading A gives a ghost copy with normal dispersion, not a tachyon in the sense of r20.  Applying F3 by
dimension to every quantity is a complex change of units and changes nothing; a reading has to state what is continued, and B continues
the kinematics.

\paragraph{Theoremlet v4.9.0-r83.4 (ghost condensation under reading B).}
Exploration (sympy).  A healthy matter side \(P_M(X)=X/2+cX^2\), \(c>0\), becomes under B \(P_T(X)=-X/2+cX^2\), the form of ghost
condensation (Arkani-Hamed, Cheng, Luty, Mukohyama, hep-th/0312099): minimum at \(X_0=1/(4c)>0\), \(P_{XX}=2c>0\), \(\rho=1/(16c)\),
\(w=-1\).  The condensate \(\langle\dot\chi\rangle\neq0\) selects a frame (candidate: the comoving frame of r81.5, where the M--T offset is
fixed in conformal time), and its fluctuations carry positive energy, so the decay of r83.1 does not occur.  By the literature
(Furukawa et al., 1001.4634) the fluctuations behave like dark matter, with \(10\,\mathrm{eV}\lesssim M\lesssim100\) GeV.  Under A the only
extremum lies at \(X<0\) with \(P_{XX}<0\).  The term \(cX^2\) is not in W1 or W2 and is an input; the height of the \(\Lambda\)-like term
is not fixed.

\paragraph{Proof and measurement sources.}
Scripts \texttt{k191}, \texttt{k191b}, \texttt{k191c}, \texttt{k192}, \texttt{k192b} and the notes in
\texttt{notes/\allowbreak research/\allowbreak v4\_10\_0\_r83/}; pre-registrations and results H76, H77 in
\texttt{notes/\allowbreak research/\allowbreak predictions/}.  \(\square\)

\subsection*{v4.9.0-r83 correction to r82}
\label{subsec:v490-r83-correction}

\paragraph{Correction to v4.9.0-r82.1(b), (c).}
r82.1(b) took the tachyon image of the Higgs sector in reading A: the whole energy--momentum tensor of the image changes sign, including
its vacuum energy, so that the vacuum energy is \((a-b)V_{\min}\) and vanishes at photon weights.  By r83.3 the image is to be formed in
reading B.  Then a potential without derivatives keeps its sign, and the homogeneous Higgs image obeys \(\ddot\phi=+V'(\phi)\): around
\(\phi=0\) it has \(\ddot\phi=-2\mu^2\phi\) and is stable (script \texttt{k192b}).  The tachyonic mass of the Higgs in M is a normal mass in
the image, the image does not condense, and the vacuum energy is \(a\,V_{\min}\), not \((a-b)V_{\min}\).  \textbf{The exact cancellation of
r82.1(b) and the transient residual of r82.1(c) are withdrawn.}  r82.1(a), r82.2 (H74), r82.3 (H75) and r82.4 stand: under B a running T
mode is per mode the negative-energy oscillator of H74, and moving T matter has the negative inertial and gravitational mass of H75.
The register line ``Higgs vacuum energy in the M--T pairing: cancels exactly at photon weights'' of r82 is superseded by the register
below.

\paragraph{Consequence.}
The structure now offers a frame and stability (ghost condensation under B), but the mechanism against the excess of about
\(10^{55}\) is gone.  A scale other than the Planck scale appears only as an upper bound from stability (\(\lesssim14\) MeV from data,
\(\lesssim2\) meV by the naturalness argument of Kaplan and Sundrum, hep-th/0505265).  OP5 is not closed.

\subsection*{v4.9.0-r83 status in the three frames}
\label{subsec:v490-r83-frames}

\begin{itemize}[leftmargin=2em]
\item r83.3 fixes how frame T enters an action: by the rapidity map of r20 on the kinematics, with constants and potentials kept.  The
photon frame remains the equal-weight mixture; nothing in this layer is evaluated there.
\item r83.1 and r83.4 work in frame M with the T sector entering through reading B; the preferred frame of r83.1(b) and r83.4 is a
frame of the cosmological background, not one of the frames M, T, P.  The bound of v4.9.0-r09 (Reading V) on a direct coupling of the
observer's rapidity to such a frame applies.
\end{itemize}

\subsection*{v4.9.0-r83 register entries}
\label{subsec:v490-r83-register}

\begin{center}
\small
\begin{tabular}{@{}>{\raggedright\arraybackslash}p{0.34\textwidth}>{\raggedright\arraybackslash}p{0.58\textwidth}@{}}
\toprule
statement & status after v4.9.0-r83\\
\midrule
quantum decay of the M--T vacuum & unbounded without a preferred frame; with one, \(c\,\Lambda^8/M_P^4\), \(\Lambda\lesssim14\) MeV (H76)\\
tachyon image of an action & reading B (kinematics continued); A excluded (H77)\\
Higgs vacuum energy in the M--T pairing & no cancellation; \(a\,V_{\min}\) (correction to r82.1)\\
ghost condensation under B & stable, frame-defining, \(w=-1\); \(cX^2\) input, height free\\
\bottomrule
\end{tabular}
\end{center}

\subsection*{v4.9.0-r83 ledger and guard update}
\label{subsec:v490-r83-ledger-guard}

\[
R_{\rm chk}^{v490\text{-}r83}=(1,1,0)
\]
for (quantum decay decided for the gravitational channel; reading of the tachyon image decided; scale not fixed).

The layer spends no guard:
\[
R_{\rm glob}^{(v490\text{-}r83,+)}=(0,1,1,1,1),
\qquad
\Gamma_{\rm guard}^{v490\text{-}r83}=(1,1,1,1,1),
\qquad
\beta_4=1/256\ \text{unspent}.
\]
It does not imply exact finite \(C5\) zero closure, exact all-readout terminality,
\(C6/T2\), a global photon-clock/RH theorem, or global all-readout closure.
